\subsubsection{Matrix Form of the Flat Interval in General Coordinates}
\label{sec:matrix-coordinate-application}
As an optional translation into ordinary matrix coordinates, the map
\(\pv{X}=f+\vct{g}\cdot\sigv\) permits the same flat interval to be
written in terms of the original four coordinate labels
\[
\vct{x}=\mcol{t,r_1,r_2,r_3},
\qquad
\dd\vct{x}=\mcol{\dd t,\dd r_1,\dd r_2,\dd r_3}.
\]
The transpose operation was defined in Eq.\ \eqref{eq:matrix-transpose}.
After \(\dd f\) and \(\dd\vct{g}\) are expanded in these four increments,
their squared difference can be written using one real symmetric matrix.
That matrix is defined by
\begin{equation}
\dd f^2-\dd\vct{g}^2
=:\mtrans{\dd\vct{x}}\Qmat(\vct{x})\,\dd\vct{x}.
\label{eq:pullback-matrix-function}
\end{equation}
It is conventionally called the pullback matrix because it expresses the
interval from the \((f,\vct{g})\) components in terms of the original
\((t,\vct{r})\) labels. In this section the name means exactly the displayed
equality \cite{doran2003}. Consequently,
\begin{equation*}
\dd s^2=\mtrans{\dd\vct{x}}\Qmat(\vct{x})\,\dd\vct{x}.
\end{equation*}
The four-by-four matrix of first derivatives of
\((f,\vct{g})\) with respect to \((t,\vct{r})\) is now considered. At a
point where this derivative matrix has an inverse, the coordinate map is
called locally invertible in the differential sense used here. The matrix
\(\Qmat\) is obtained by placing this derivative matrix on both sides of the
one-positive, three-negative interval matrix, with its transpose on the
left. An invertible real change of coordinates therefore cannot change the
count of one positive and three negative squared directions. That sign count
is called the signature \((+,-,-,-)\). The entries of \(\Qmat\) may
nevertheless depend on position when nonlinear, or curvilinear, coordinates
are used. Such position dependence alone does not imply curvature because
\((f,\vct{g})\) still supplies components in the fixed flat algebra. If the
derivative matrix has no inverse, a nonzero coordinate direction is sent to
zero. The same direction is then sent to zero by \(\Qmat\), so
\(\Qmat\) necessarily has no inverse. This case is called degenerate.

For a real symmetric matrix, four mutually perpendicular directions can be
chosen so that multiplication by \(\Qmat\) only rescales each direction.
The rescaling numbers are called eigenvalues, and the corresponding
directions are called eigenvectors. At every locally invertible point, the
ordered eigenvalue function is defined by
\begin{equation}
\lambda(n):=\eqtext{ordered eigenvalue }n\eqtext{ of }\Qmat,
\qquad n\in\{1,2,3,4\}.
\label{eq:metric-eigenvalue-function}
\end{equation}
For the four-entry use required here, the diagonal matrix function is
defined by
\begin{equation}
\eqfunc{diag}(A,B,C,D):=
\begin{bmatrix}
A&0&0&0\\
0&B&0&0\\
0&0&C&0\\
0&0&0&D
\end{bmatrix}.
\label{eq:matrix-diag-four}
\end{equation}
The normalized eigenvectors are placed in the columns of \(\Nmat\), giving
\[
\Qmat=\Nmat\eqfunc{diag}(\lambda(1),\lambda(2),\lambda(3),\lambda(4))\,\mtrans{\Nmat},
\qquad
\mtrans{\Nmat}\Nmat=\I.
\]
The second equality states that the columns have unit length and are
mutually perpendicular; \(\I\) is the identity matrix of the required
four-by-four size. The eigenvalues are ordered numerically so that
\[
\lambda(1)>0>\lambda(2)\geq\lambda(3)\geq\lambda(4).
\]
The corresponding direction function is defined by
\begin{equation}
\vct{u}(n):=\eqtext{column }n\eqtext{ of }\Nmat.
\label{eq:metric-direction-function}
\end{equation}
The defining eigenvalue relation is therefore
\[
\Qmat\vct{u}(n)=\lambda(n)\vct{u}(n).
\]
If two or more eigenvalues are equal, the individual eigenvectors within
their shared directions are not unique. Any mutually perpendicular unit
vectors spanning those same directions give the same matrix and interval.
The local differential directions are defined by
\begin{equation}
w(n):=\sqrt{|\lambda(n)|}\,\mtrans{\vct{u}(n)}\dd\vct{x}.
\label{eq:local-differential-directions}
\end{equation}
Each \(w(n)\) is the coordinate increment along one eigenvector,
rescaled by the square root of the corresponding eigenvalue magnitude.
Substitution of the matrix decomposition gives
\begin{align*}
\dd s^2
&=\lambda(1)\bigl(\mtrans{\vct{u}(1)}\dd\vct{x}\bigr)^2
+\lambda(2)\bigl(\mtrans{\vct{u}(2)}\dd\vct{x}\bigr)^2\notag\\
&\quad+\lambda(3)\bigl(\mtrans{\vct{u}(3)}\dd\vct{x}\bigr)^2
+\lambda(4)\bigl(\mtrans{\vct{u}(4)}\dd\vct{x}\bigr)^2\notag\\
&=w(1)^2-w(2)^2-w(3)^2-w(4)^2.
\end{align*}
An interval with one positive and three negative squared directions is
called Lorentzian. A Euclidean squared length would instead contain only
plus signs. The three minus signs therefore cannot be dropped by changing
matrix directions. When \(\Qmat\) varies with position, the four quantities
\(w(n)\) describe the interval at one point; a single set of
coordinates whose differentials equal all four \(w(n)\) need not exist
throughout an extended region.
